% The manuscript, in the layout you write and review it in: plain article, figures and tables next to the text that
% cites them. It adds no content of its own, it reads sections/ and latex/, and it is the only document here until a
% journal is chosen.
%
% When one is: copy the matching file out of templates/documents/ into this directory and build both. That file
% supplies the class, the bibliography style and the renderers for \projauthor, \projauthorcorr, \projtable and
% \projdeclaration; nothing in sections/ changes, and logs/build.sh and the tests pick the new document up on their
% own. A submission class is not a reading layout: cas-sc, for one, empties the float specifier of figures so they
% collect at the end, which is what a publisher wants and a reader does not. So this file stays either way.
\documentclass[11pt,a4paper]{article}

\usepackage[T1]{fontenc}
\usepackage[utf8]{inputenc}
\usepackage{lmodern}
\usepackage{amsmath,amssymb}
\usepackage{graphicx}
\usepackage{array}
\usepackage{tabularx}
\usepackage{booktabs}
\usepackage{longtable}
\usepackage{multirow}
\usepackage{float}
\usepackage{xcolor}
\usepackage[margin=2.4cm]{geometry}
\usepackage[hidelinks]{hyperref}
\usepackage[authoryear]{natbib}
\usepackage{caption}

% Let a float sit near its citation: the default limits reserve a fifth of every page for text, which on a
% figure-dense manuscript pushes floats forward until they pile up at the end.
\setcounter{topnumber}{3}
\setcounter{totalnumber}{4}
\renewcommand{\topfraction}{0.9}
\renewcommand{\bottomfraction}{0.8}
\renewcommand{\textfraction}{0.07}
\renewcommand{\floatpagefraction}{0.8}

\captionsetup{font=small, labelfont=bf}
\setlength{\parskip}{0.3em}

% The front matter as data: title, authors, their e-mail, ORCID and contributions, and the one affiliation.
% No journal's syntax appears here, because every class asks for the same facts in a different shape: one wants
% a single name per author with the address as key=value pairs, another wants the given and family names split
% into separate commands and the address broken into fields of its own, and the reading copy wants one
% \and-separated line. So this file states each fact once and a template renders it, exactly as the abstract and
% the keywords already work. Supporting one more journal means writing the renderers below in the new
% template, and changing nothing here.
%
% A template must define, before it uses \projauthorlist:
%   \projauthor{initials}{given}{family}{email}{orcid}{CRediT roles}      one non-corresponding author
%   \projauthorcorr{...}                                                 same six, for the corresponding author
% ORCID may be empty, so a renderer that prints it guards the field with \projifhasorcid.
%
% THIS IS THE FIRST FILE TO EDIT IN A NEW PAPER. Nothing else holds any of these facts.

\newcommand{\projtitle}{The Title of the Paper, Stated in Full and Without a Subtitle Unless One Earns Its Place}
\newcommand{\projshorttitle}{Short Running Title}
\newcommand{\projrunningauthors}{Family-Name et al., 2026}

\newcommand{\projaffildiv}{Department or Institute}
\newcommand{\projaffilorg}{University}
\newcommand{\projaffilstreet}{Street and Number}
\newcommand{\projaffilcity}{City}
\newcommand{\projaffilstate}{Region}
\newcommand{\projaffilpostcode}{00000}
\newcommand{\projaffilcountry}{Country}

% Not every author has an ORCID on file, and a renderer that prints the field unconditionally would emit an empty
% link for the ones who do not. This runs its second argument only when the first is a non-empty ORCID.
\newcommand{\projifhasorcid}[2]{\def\projorcidarg{#1}\ifx\projorcidarg\empty\else#2\fi}

% The authors in the order they appear on the paper. The corresponding author is marked by which macro the line
% calls, not by a flag, so no layout has to decide what "corresponding" looks like in its own notation.
% The roles are the CRediT taxonomy's own terms: they go into the author block of a class that has one, and into
% the contributions declaration, which is rendered from this same list.
\newcommand{\projauthorlist}{%
	\projauthorcorr{A.B.}{Ada}{Firstauthor}{ada.firstauthor@example.org}{0000-0000-0000-0000}%
		{Conceptualization, Methodology, Software, Formal analysis, Writing -- original draft}%
	\projauthor{C.D.}{Carl}{Secondauthor}{carl.secondauthor@example.org}{}%
		{Validation, Writing -- review \& editing, Supervision}%
}


% The front matter, from sections/manuscript_metadata.tex. The reading copy prints the names on one line, so
% its renderer emits only those and separates them with \and from the second author on; the e-mail, the ORCID
% and the CRediT roles belong to the submission layouts and to the contributions declaration.
\newcommand{\projsep}{}
\newcommand{\projauthor}[6]{\projsep #2 #3\renewcommand{\projsep}{\and}}
\newcommand{\projauthorcorr}[6]{\projauthor{#1}{#2}{#3}{#4}{#5}{#6}}

\title{\textbf{\projtitle}}
\author{\projauthorlist}
\date{\projaffildiv,\\
	\projaffilorg, \projaffilcity, \projaffilcountry\\[0.6em]
	\textit{Reading copy, figures in place. Not the submission layout.}\\
	\today}

% A generated table, typeset at the full text width. Every table in latex/ is a tabularx or a tabular* that
% already fills \textwidth, so no template has to box or scale one; this hook stays so that a journal needing
% something of its own can redefine it in one place.
\newcommand{\projtable}[1]{\input{#1}}

% One declaration: the statement is shared, the heading belongs to the layout. The declarations are
% listed in sections/manuscript_declarations.tex and this is the only place a layout decides how they
% are presented.
\newcommand{\projdeclaration}[2]{\section*{#1}\input{#2}}

% Every macro file the pipeline writes, read once so that no document can be missing one. A new step that emits
% numbers adds its file here and nowhere else; tests/test_manuscript_layout.py fails if latex/ holds a
% *_variables.tex that this list does not name.
% Written by scripts/99_demo_variables.py. Do not edit: the next run overwrites it.
\newcommand{\DemoHeadline}{0.74}
\newcommand{\DemoNSamples}{859}



\begin{document}
\maketitle

\begin{abstract}
% The abstract text only, with no environment of its own. Some classes want it wrapped in an environment and
% others in a command, and a class that typesets the environment where it stands would put the abstract before
% the title, so the choice belongs to the template.
%
% Write it last. One useful order: the problem, what was done, three results, what the reader may not conclude.
% Every number here is a macro from latex/, never a typed digit: \DemoHeadline is an example of one.
The problem this paper addresses, in one sentence that a reader outside the field can follow. What was done, and
on what data. The main result, which came out at \DemoHeadline\ on \DemoNSamples\ samples. The second result. The
third. What the reader should not conclude from any of it.


\medskip\noindent\textbf{Keywords:} % The keyword list only, comma separated, for whatever the template's keyword command or environment is.
first keyword, second keyword, third keyword, fourth keyword, fifth keyword

\end{abstract}

% The manuscript itself, independent of any journal template: the sections in order, each one free of preamble.
% A template supplies the class, the author block and the bibliography style; this file supplies the content.
% The declarations (funding, competing interests and the rest) are not here: a layout wraps them in its own
% way, so it reads sections/manuscript_declarations.tex separately, after this file.
%
% This list is the order of the paper. tests/test_manuscript_layout.py checks that the numbered files are listed
% in ascending order, so renumber the files when you reorder them.
% The Introduction, with no preamble of its own: a template supplies the class, the front matter and the bibliography
% style. Cite with \citep{} and \citet{}; refer to a generated table with \projtable, never \input.
\section{Introduction}

Replace this paragraph. Every number in the text comes from a macro that a script in scripts/ wrote into latex/;
the check step fails the build if the text names a macro that nothing produced. Cite as
\citet{example2026} or \citep{example2026}, from references.bib.

% The Methods, with no preamble of its own: a template supplies the class, the front matter and the bibliography
% style. Cite with \citep{} and \citet{}; refer to a generated table with \projtable, never \input.
\section{Methods}

Replace this paragraph. Every number in the text comes from a macro that a script in scripts/ wrote into latex/;
the check step fails the build if the text names a macro that nothing produced.

% The Results, with no preamble of its own: a template supplies the class, the front matter and the bibliography
% style. Cite with \citep{} and \citet{}; refer to a generated table with \projtable, never \input.
\section{Results}

Replace this paragraph. Every number in the text comes from a macro that a script in scripts/ wrote into latex/;
the check step fails the build if the text names a macro that nothing produced.

\begin{table}[htbp]
	\centering
	\caption{An example of a generated table. The caption is written by hand here, next to the item it describes,
		while the body comes from latex/ and is overwritten on every run; a caption that states a number states it
		with a macro, as this one does: \DemoNSamples{} samples.}
	\label{tab:demo}
	\projtable{latex/demo_table.tex}
\end{table}

% The Discussion, with no preamble of its own: a template supplies the class, the front matter and the bibliography
% style. Cite with \citep{} and \citet{}; refer to a generated table with \projtable, never \input.
\section{Discussion}

Replace this paragraph. Every number in the text comes from a macro that a script in scripts/ wrote into latex/;
the check step fails the build if the text names a macro that nothing produced.

% The Conclusions, with no preamble of its own: a template supplies the class, the front matter and the bibliography
% style. Cite with \citep{} and \citet{}; refer to a generated table with \projtable, never \input.
\section{Conclusions}

Replace this paragraph. Every number in the text comes from a macro that a script in scripts/ wrote into latex/;
the check step fails the build if the text names a macro that nothing produced.



% The declarations publishers ask for, in the order they are most often requested. Each one is a file of its own so
% that a statement is written once and every layout reports the same thing; the heading, and whether the set is
% gathered under one of its own, is the layout's business, which is what \projdeclaration supplies. Adding a
% journal means defining that command and reordering these lines, not rewriting the statements.
\projdeclaration{Availability of data and materials}{sections/06_1_Data_Available.tex}
\projdeclaration{Competing interests}{sections/06_2_Competing_Interests.tex}
\projdeclaration{Funding}{sections/06_3_Funding.tex}
\projdeclaration{Authors' contributions}{sections/06_4_Authors_Contributions.tex}
\projdeclaration{Acknowledgements}{sections/06_5_Acknowledgements.tex}
\projdeclaration{Declaration of generative AI use}{sections/06_6_Generative_AI.tex}


\bibliographystyle{plainnat}
\bibliography{references.bib}

\end{document}
